\documentclass[10pt,a4paper]{article}
\usepackage[margin=18mm]{geometry}
\usepackage{amsmath,amssymb,booktabs,array}
\usepackage{cite}
\renewcommand{\citepunct}{], [}
\renewcommand{\citedash}{]--[}
\input{glyphtounicode}
\pdfgentounicode=1
\usepackage[hidelinks]{hyperref}
\hypersetup{pdftitle={An elementary VSL lane-drop task: model and sensitivity},
 pdfauthor={Differentiable traffic task review},pdfsubject={Compact methodology}}
\setlength{\parindent}{0pt}
\setlength{\parskip}{3pt}
\setlength{\abovedisplayskip}{4pt}
\setlength{\belowdisplayskip}{4pt}
\setlength{\abovedisplayshortskip}{3pt}
\setlength{\belowdisplayshortskip}{3pt}
\newcommand{\pos}[1]{[#1]_+}
\newcommand{\bud}{\mathcal B}
\newcommand{\dd}{\,\mathrm d}
\newcommand{\head}[1]{\par\vspace{4pt}\noindent\textbf{#1}\par\nobreak\vspace{1pt}}
\begin{document}
\begin{center}
{\large\bfseries An elementary VSL lane-drop task: model and sensitivity}
\end{center}
\vspace{-5pt}

\head{1. Research question}
Can the distribution of traffic between continuing lanes change useful
common variable-speed-limit (VSL) advice, even when ending-lane traffic and
total continuing traffic agree? The task connects lane information to
receiving, merge service and complete travel time, then differentiates that
same causal chain for control.

\head{2. Task and global assumptions}
A permanent three-to-two drop has ending lane E, adjacent continuing lane A
and far continuing lane B. Preparation ends at $L_p$; a terminal of length
$b$ ends at the choke $L=L_p+b$. At time zero, independent lane densities
$(e_0,a_0,b_0)\geq0$ are uniform on the common footprint $[0,\ell]$,
$0<\ell\leq L_p$. Set $r=e_0+a_0$ and total mass
$M=\ell(e_0+a_0+b_0)>0$. The terminal and recovery history start empty;
there are no later external arrivals. Choose one constant common advice
$v\in[v_-,v_+]$, $0<v_-<v_+$, with fixed initial state, geometry and physical parameters.
The objective counts every initial unit's time until it crosses the choke.

\textbf{Global assumptions.}
Traffic is continuous mass; the remaining preparation body is homogeneous.
Unblocked motion follows the common advice, with fixed backward-wave
parameter $w$ and jam density $k_j$. Only E-to-A transfer occurs, before
the terminal and at the choke; B has separate service. Point service uses
current-state priority/cooperation and representative free/queued speed
marks. Forced E entries alone create recovery loss, retained as age/mark
history behind one shared service budget.

Triangular-diagram and cumulative-boundary ingredients follow the link
transmission model (LTM)~\cite{jin}; the advice-response family is selected
here, with algebraic capacity/inverse maps:
\begin{equation}
 K=wk_j,\quad k_c(v)=\frac{K}{w+v},\quad C(v)=vk_c(v),\qquad
 c(u)=\frac{Ku}{w+u},\quad U(g)=\frac{wg}{K-g}.
 \label{eq:fd}
\end{equation}
$R$ and $R_B$ are undisturbed A/E and B service ceilings, respectively.
Use positive physical scales and the sufficient operating domain
\begin{equation}
 e_0,a_0,b_0\leq k_c(v_+),\qquad
 v_+r<R\leq C(v_-),\qquad v_+b_0\leq R_B\leq C(v_-),\qquad M<k_jb.
 \label{eq:domain}
\end{equation}
It keeps preparation free, terminal receiving available and B unqueued.
Work and recovery choices are specified below; $[x]_+=\max(x,0)$.

\head{3. Traffic model}
Preparation determines lane-specific entries; terminal travel determines
eligibility; forced service and recovery then update each other causally.

\textbf{3.1. Pre-choke merging.}
Let $e(t),a(t)$ be the remaining preparation body's E/A densities.
Own A flow uses receiving first. E transfer is limited by its donor offer
and by remaining adjacent service:
\begin{equation}
 S(v,a)=\min\{C(v),w(k_j-a)\},\qquad
 f=\min\!\left\{ve,\frac{S(v,a)-va}{\gamma}\right\},\qquad
 \dot e=-f/h,\quad \dot a=f/h.
 \label{eq:prep}
\end{equation}
Here $h>0$ is a physical interaction length, $\gamma\geq1$ is a receiving-work
factor, and $f/h$ is the lateral density source. Thus $e+a=r$ and $b_0$
stays constant. Donor-limited decay is exponential per traveled distance;
advice affects preparation through normalized receiving headroom $k_c(v)$.
\emph{Adapted:} paired conservative lane sources follow the structure
in~\cite{subraveti}; receiving work and own-lane priority draw on
Carey et al.~\cite{carey}. The homogeneous donor/headroom minimum and
interaction length are defined here.

\textbf{3.2. Finite terminal and the merging point.}
For $j\in\{E,A,B\}$, let $F_j,P_j,Y_j$ be cumulative terminal entry,
service-eligible mass and departure. During
$(L_p-\ell)/v<t<L_p/v$,
\begin{equation}
 \dot F_E=ve(t),\quad \dot F_A=v[r-e(t)],\quad \dot F_B=vb_0;
 \qquad P_j(t)=F_j(t-b/v).
 \label{eq:arrivals}
\end{equation}
Entry rates vanish outside that interval; negative-time counts are zero.
Define physical terminal stock $N_j=F_j-Y_j$, eligible waiting mass
$q_j=P_j-Y_j$, offer $\lambda_j=\dot P_j$ and discharge $g_j=\dot Y_j$.
Departures satisfy $0\leq Y_j\leq P_j\leq F_j$, with
$\dot q_j=\lambda_j-g_j$. Domain~\eqref{eq:domain} protects finite storage
and gives $Y_B=P_B$. Early-merged E becomes ordinary A; residual E travels
through the terminal and merges at the choke.
\emph{Reused/adapted:} boundary travel and storage follow LTM~\cite{jin};
the preparation cutoff and terminal handoff are defined here.

\newpage
\textbf{3.3. Forced merging and service sharing.}
Use free marks $u_j=v$ and queued marks $u_j=U(g_j)$ for $j=A,E$.
A costs one work unit per mass; E's maneuver cost increases with adjacent
flow and a positive speed difference:
\begin{equation}
 c_E=c_0+\chi g_A\pos{u_A-u_E},\qquad
 g_A+c_Eg_E\leq\bud,\qquad
 H(g;g_A,u_A)=g\{c_0+\chi g_A\pos{u_A-U(g)}\}.
 \label{eq:work}
\end{equation}
$c_0\geq1$ is baseline E work and $\beta\geq0$ sets mismatch strength.
Set $\chi=\beta/[C(v_+)v_+]$ once and hold it fixed.
Require $c_0>\beta(1+v_+/w)$ so each feasible root $H(g)=x$
is unique. Here $x$ is work allocated to E. A fitting free A offer has
priority; queued A cooperates through equal work and unused-work redistribution:
\begin{center}
\setlength{\tabcolsep}{5pt}
\begin{tabular}{@{}ll@{}}
\toprule
A, E status & Selected service and free-admission condition\\
\midrule
Free, free & $g_A=\lambda_A,\ g_E=\lambda_E$ if
$\lambda_A+c_0\lambda_E\leq\bud$.\\
Free, queued & $g_A=\lambda_A\leq\bud,\quad
H(g_E;\lambda_A,v)=\bud-\lambda_A$.\\
Queued, queued & $g_A=\bud/2,\quad
H(g_E;\bud/2,U(\bud/2))=\bud/2$.\\
Queued, free & $g_E=\lambda_E,\ g_A=\bud-c_0\lambda_E$ if
$c_0\lambda_E\leq\bud/2$.\\
\bottomrule
\end{tabular}
\end{center}
A strict free-service shortfall creates a queue. At an actual zero hit,
retest free admission; equality retains feasible free continuation. Use
right-hand arrivals at jumps and retain recovery history across transitions.
\emph{Defined here:} the mismatch formula, representative marks, sharing
table and transition guards. Receiving-work accounting draws on~\cite{carey};
the particular state-dependent node is specified here.

\textbf{3.4. Recovery and the shared budget.}
Point-merge models link slow insertion and bounded acceleration to lost
discharge~\cite{leclercq2011,leclercq}. Differentiating the isolated passing
curve in~\cite{leclercq} gives an equivalent boundary-speed quantity;
we cap its service at $R$. For an E entry of mark $u$, age $\theta\geq0$
and recovery acceleration $a_r>0$, define
\begin{equation}
 V(\theta;u)=\sqrt{(w+u)^2+2wa_r\theta}-w,\qquad
 D(\theta;u)=\pos{R-c(V(\theta;u))}.
 \label{eq:profile}
\end{equation}
$V$ is an equivalent boundary-speed quantity; $D$ is remaining lost service.
Let $T(\xi)$ be the admission time of E mass rank $\xi$ and fix
$m=1$ vehicle-equivalent. The recovery reading and available budget are
\begin{equation}
 p_{\rm rec}(t)=\frac1m\int_{\pos{Y_E(t)-m}}^{Y_E(t)}
 D\bigl(t-T(\xi);u_E(T(\xi))\bigr)\dd\xi,\qquad
 \bud(t)=\pos{R-p_{\rm rec}(t)}.
 \label{eq:recovery}
\end{equation}
All retained profiles age; after the window fills, new E replaces oldest
mass weight. The fixed denominator also applies before fill. A consumes the
shared budget while E supplies the profiles; their predicted age/mark history
updates the account. \emph{Derived/adapted:} the isolated profile from~\cite{leclercq}.
\emph{Defined here:} capping, the E-only source and recent-mass renewal.

\head{4. How the task tests information}
Compare the same initial E mass and continuing total with different A/B
splits. A busier A can restrict preparation, leaving more E for forced work
and recovery. Advice trades these effects against ordinary travel time.
With a fixed horizon $H_*$ beyond cohort clearance, minimize complete TST:
\begin{equation}
 J(v)=\int_0^{H_*}\!\left[M-\sum_jY_j(t)\right]\dd t
 =\frac{M(L-\ell/2)}v+\int_0^{H_*}(q_A+q_E)\dd t.
 \label{eq:objective}
\end{equation}
The predictors share geometry and physical parameters and may use all
inference from admitted predecision histories: aggregate density, flow,
speed and cumulative counts, including any permitted E-resolved history.
One knows initial E and combined A+B mass; the other knows each lane.
An equal-information compact peer adds just initial A mass to the former,
reconstructing all three initial densities and the same dynamics.

\newpage
\head{5. Complete sensitivity}
The model-specific sensitivities are derived here at fixed initial traffic
and physical parameters. Carry preparation, moving entry/eligibility clocks,
service/marks, recovery aging/retirement and queue-event times through the
full objective. A subscript $v$ denotes advice sensitivity. On smooth branches,
\begin{equation}
 g_{E,v}=\frac{x_v-H_{g_A}g_{A,v}-H_{u_A}u_{A,v}}{H_g},\qquad H_g>0.
 \label{eq:rootgradient}
\end{equation}
For a filled window with $g_E(s_m)>0$, let $Y_E(s_m)=Y_E(t)-m$ and $Z_E=Y_{E,v}$.
The cutoff relation below is derived here. Differentiate~\eqref{eq:recovery}
including marks, flows, retirement and source events. For an isolated
transverse guard $G$ along the history-aware trajectory, the remaining
event corrections for continuous traffic coordinates $X$ follow~\cite{kong}:
\begin{equation}
 \begin{gathered}
 g_E(s_m)s_{m,v}=Z_E(t)-Z_E(s_m),\qquad \tau_v=-G_v^-/G_t^-,\\
 X_v^+=X_v^-+(\dot X^- -\dot X^+)\tau_v.
 \end{gathered}
 \label{eq:eventgradient}
\end{equation}
Thus the complete derivative, on a fixed regular event/history sequence, is
\begin{equation}
 J'(v)=-\int_0^{H_*}\!\sum_jY_{j,v}(t)\dd t
 =-\frac{M(L-\ell/2)}{v^2}+\int_0^{H_*}(q_{A,v}+q_{E,v})\dd t.
 \label{eq:objectivegradient}
\end{equation}
\makeatletter
\renewcommand{\section}{\@ifstar{\head}{\head}}
\makeatother
% The reference list is excluded from the main-methodology page budget.
\begin{thebibliography}{6}
\setlength{\itemsep}{0pt}\setlength{\parskip}{0pt}
\bibitem{jin} W.-L. Jin, ``Continuous formulations and analytical properties
of the link transmission model,'' \emph{Transp. Res. B, Methodol.}, vol.~74,
pp.~88--103, 2015, doi: \href{https://doi.org/10.1016/j.trb.2014.12.006}{10.1016/j.trb.2014.12.006}.
\bibitem{subraveti} H.~H.~S. Nagalur Subraveti, V.~L. Knoop, and B. van Arem,
``First order multi-lane traffic flow model---an incentive based macroscopic
model to represent lane change dynamics,'' \emph{Transportmetrica B, Transp.
Dyn.}, vol.~7, no.~1, pp.~1758--1779, 2019, doi:
\href{https://doi.org/10.1080/21680566.2019.1700846}{10.1080/21680566.2019.1700846}.
\bibitem{carey} M. Carey, C. Balijepalli, and D. Watling, ``Extending the cell
transmission model to multiple lanes and lane-changing,'' \emph{Netw. Spat.
Econ.}, vol.~15, pp.~507--535, 2015, doi:
\href{https://doi.org/10.1007/s11067-013-9193-7}{10.1007/s11067-013-9193-7}.
\bibitem{leclercq2011} L. Leclercq, J.~A. Laval, and N. Chiabaut,
``Capacity drops at merges: An endogenous model,'' \emph{Procedia Soc.
Behav. Sci.}, vol.~17, pp.~12--26, 2011, doi:
\href{https://doi.org/10.1016/j.sbspro.2011.04.505}{10.1016/j.sbspro.2011.04.505}.
\bibitem{leclercq} L. Leclercq, V.~L. Knoop, F. Marczak, and S.~P. Hoogendoorn,
``Capacity drops at merges: New analytical investigations,'' \emph{Transp.
Res. C, Emerg. Technol.}, vol.~62, pp.~171--181, 2016, doi:
\href{https://doi.org/10.1016/j.trc.2015.06.025}{10.1016/j.trc.2015.06.025}.
\bibitem{kong} N.~J. Kong, J.~J. Payne, J. Zhu, and A.~M. Johnson,
``Saltation matrices: The essential tool for linearizing hybrid dynamical
systems,'' \emph{Proc. IEEE}, vol.~112, no.~6, pp.~585--608, Jun. 2024, doi:
\href{https://doi.org/10.1109/JPROC.2024.3440211}{10.1109/JPROC.2024.3440211}.
\end{thebibliography}
\end{document}
